\documentclass[a4paper,11pt]{article}
\usepackage[T1]{fontenc}
\usepackage[utf8]{inputenc}
\usepackage{t1enc}
\usepackage[dvips]{color}
\usepackage{amsmath,amssymb,makeidx,stmaryrd}
\usepackage[normalem]{ulem}
\usepackage{fancybox}
\definecolor{gris}{gray}{0.9}
\def\bas{\fcolorbox{gris}{gris}{\parbox{0.986\textwidth}{
	\color{red}\footnotesize\it mardi 18 mars 1997\hfill Page : \thepage/1}}}
\def\haut{\fcolorbox{gris}{gris}{\parbox{0.986\textwidth}{
	\color{red}
	\bf Concours Général de Mathématiques \hfill mardi 18 mars 1997
	}}}
\makeatletter
\def\ps@cg{%
\let\@oddhead\haut%
\let\@evenhead\haut%
\let\@oddfoot\bas
\let\@evenfoot\bas}
\makeatother
\pagestyle{cg}
\usepackage[left=3.5cm, right=3.5cm, top=3cm, bottom=3cm]{geometry}
\newcommand{\vect}[1]{\overrightarrow{\,\mathstrut#1\,}}
\renewcommand{\theenumi}{\textbf{\arabic{enumi}}}
\renewcommand{\labelenumi}{\textbf{\theenumi.}}
\renewcommand{\theenumii}{\textbf{\alph{enumii}}}
\renewcommand{\labelenumii}{\textbf{\theenumii.}}
\parindent0pt
\everymath{\displaystyle}
\usepackage[french]{babel}

\DecimalMathComma
\usepackage[np]{numprint}
\begin{document}

\begin{center}\textbf{{\color{blue} Exercice I}}
\end{center}

On a placé un jeton sur chaque sommet d'un polygone régulier à \np{1997} côtés. 

Sur chacun de ces jetons est inscrit un entier relatif, la somme de ces entiers relatifs étant égale à 1. On choisit un sommet de départ et on parcourt le polygone dans le sens trigonométrique en ramassant les jetons au fur et à mesure tant que la somme des entiers inscrits sur les jetons ramassés est strictement positive.

Peut-on choisir le sommet de départ de façon à ramasser tous les jetons ? Si oui, combien y-a-t-il de choix possibles ?

\begin{center}\textbf{{\color{blue} Exercice II}}
\end{center}

Une capsule spatiale a la forme du solide de révolution délimité par une sphère de centre 0, de rayon R, et un cône de sommet O qui rencontre cette sphère selon un cercle de rayon $r$.

Quel est le volume maximal d'un cylindre droit contenu dans cette capsule, le cylindre et la capsule ayant même axe de révolution ?

\begin{center}\textbf{{\color{blue} Exercice III}}
\end{center}

C est un cube d'arête 1 et $p$ est la projection orthogonale sur un plan. Quelle est la
valeur maximale de l'aire de $p$(C) ?

\begin{center}\textbf{{\color{blue} Exercice IV}}
\end{center}

Etant donné un triangle ABC, on note $a, b, c$ les longueurs de ses côtés et $m, n, p$ les longueurs de ses médianes. 

Pour tout réel $\alpha$ strictement positif, on définit le réel $\lambda(\alpha)$ par la relation :

\[a^{\alpha} + b^{\alpha} + c^{\alpha} = [\lambda(\alpha)]^{}\left(m^{\alpha} + n^{\alpha} + p^{\alpha} \right)\]

\begin{enumerate}
\item Calculer $\lambda(2)$.
\item Calculer la limite de $\lambda(\alpha)$ lorsque $\alpha$ tend vers $0$.
\item A quelle condition portant sur $a, b, c$ le réel $\lambda(\alpha)$ est-il indépendant de $\alpha$ ?
\end{enumerate}

\begin{center}\textbf{{\color{blue} Exercice V}}
\end{center}

Dans le plan, soient A et B deux points distincts.

Pour tout point C extérieur à la droite (AB), on note G l'isobarycentre du triangle ABC et I le centre de son cercle inscrit.

\medskip

\begin{enumerate}
\item Soit $\alpha$ un réel tel que $0 < \alpha < \pi$. 

Quel est l'ensemble $\Gamma$ des points C tels que

\[\left(\vect{\text{CA}},\:\vect{\text{CB}}\right) = \alpha +2k\pi.\]

$k$ étant un entier ? 

Lorsque C décrit $\Gamma$, montrer que G et I décrivent deux arcs de
cercle que l'on précisera.
\item On suppose désormais que $\dfrac{\pi}{3} < \alpha < \pi$. Comment doit-on choisir C dans $\Gamma$pour que la distance GI soit minimale ?
\item On note $f(\alpha)$ la distance minimale GI de la question précédente. 

Expliciter $f(\alpha)$ en fonction de $a =$ AB et $\alpha$. 

Déterminer la valeur maximale de $f(\alpha)$ lorsque $\alpha$ décrit $\left]\dfrac{\pi}{3}~;~\pi\right[$.
\end{enumerate}
\end{document}